% !TeX root = ../main.tex

\chapter{System Design and Implementation}\label{chapter:system-design}

This chapter documents the system that operationalises the configurations of Chapter~3.
The objective is reproducibility: the architectures are described at a level of detail
sufficient for an independent team to re-implement them, and the supporting
infrastructure (prompt registry, model client, observability layer, dashboard, and test
suite) is described at the same level. Code references in this chapter point to the
thesis artifact released alongside the document.

\section{Architectural Patterns Evaluated}\label{sec:sys:patterns}

The five configurations populate three orthogonal axes of the agentic design space:
\emph{number of stages} (single vs.\ multi-stage), \emph{routing} (none, static, or
LLM-driven), and \emph{specialisation} (a single generalist agent vs.\ several domain
specialists vs.\ a consolidated specialist prompt). Table~\ref{tab:configs} summarises
the five configurations and their position in this space.

\begin{table}[h]
  \centering
  \small
  \begin{tabular}{lllll}
    \toprule
    \textbf{ID} & \textbf{Name} & \textbf{Stages} & \textbf{Routing} &
    \textbf{Specialisation} \\
    \midrule
    B1 & Zero-shot single agent            & 1 & none          & generalist \\
    B4 & Chain-of-thought single agent      & 1 & none          & generalist + reasoning \\
    M0 & Classify-then-extract pipeline    & 2 & static (gate) & generalist \\
    M1 & Orchestrator + 3 specialists      & 2 & LLM router    & three domain specialists \\
    M6 & Consolidated specialist (ablation) & 1 & none          & one agent, union of specialist prompts \\
    \bottomrule
  \end{tabular}
  \caption{The five evaluated configurations across the three design axes.}
  \label{tab:configs}
\end{table}

\section{Baseline Architectures}\label{sec:sys:baselines}

\paragraph{B1 --- ContractEval replication.}
B1 is a faithful replication of the zero-shot prompt protocol of
\cite{liu2025contracteval}. The model receives a single user message containing the
contract text, the target category name, the category indicator (a short
attorney-style description of what the category covers, taken from CUAD's annotation
guidelines), and an explicit instruction to either return the verbatim span(s) of the
clause or the literal string ``no related clause''. The system message is identical
across categories. No examples, no chain-of-thought instructions, and no auxiliary
structure are added. B1 establishes the floor for a competent single-agent reading.

\paragraph{B4 --- Chain-of-thought single agent.}
B4 retains B1's task framing but explicitly elicits step-by-step reasoning followed by
a structured final answer. The user prompt instructs the model to reason about the
category definition, scan the contract for candidate spans, weigh ambiguous matches, and
only then emit a final \texttt{<answer>...</answer>} block containing either verbatim
span(s) or ``no related clause''.

The B4 prompt went through one substantive revision during pilot work. The original
delimiter-based parser extracted any text occurring after a sentinel string, which under
certain models produced spurious matches: when the chain-of-thought monologue itself
quoted contract passages, the parser harvested them as if they were the final answer,
inflating the over-extraction rate. The fix was to enforce a structured-output contract:
the parser now requires the final block to be enclosed in an explicit answer tag and
rejects the trace otherwise, with a small repair step that re-prompts the model for a
structured emission if the tag is missing. Only the post-fix version is in the official
runs; the parsing-bug runs are retained for reference but excluded from all reported
statistics.

\section{Multi-Agent Architectures}\label{sec:sys:multiagent}

\paragraph{M0 --- Classify-then-extract pipeline.}
M0 decomposes the per-category task into two LLM calls. The first call is a binary
\emph{classifier} agent that receives the contract and the category and returns a single
yes/no token indicating whether any span of the category is present. The second call is
an \emph{extractor} agent that is invoked only when the classifier returns yes; it
receives the same inputs and produces the verbatim span(s). When the classifier returns
no, the extractor is bypassed and the system records ``no related clause'' as the final
answer. The two agents share a category indicator but have disjoint prompts: the
classifier's prompt explicitly forbids span emission, and the extractor's prompt assumes
presence. M0 implements the simplest form of conditional computation in the design space
and is the architectural reference for the value of a classification gate.

\paragraph{M1 --- Orchestrator with three specialists.}
M1 consists of an \emph{orchestrator} that receives the contract and the category and
routes the request, via an LLM router prompt, to one of three domain specialists. The
router prompt enumerates the three specialists, gives a one-line description of each,
and asks for a single category-name decision. The three specialists partition the 41
CUAD categories as follows:

\begin{itemize}
  \item \textbf{Risk/Liability specialist} (13 categories): liability caps,
    indemnification, insurance, warranty, governing law, and similar risk-allocation
    provisions.
  \item \textbf{Temporal/Renewal specialist} (11 categories): term, renewal,
    termination, notice periods, post-termination obligations, and other time-bounded
    provisions.
  \item \textbf{IP/Commercial specialist} (17 categories): IP ownership and licence
    grants, exclusivity, non-compete, audit rights, and commercial covenants.
\end{itemize}

Categories were assigned to specialists by the author in consultation with the CUAD
category descriptions; the assignment was fixed before any model was run. Each
specialist's prompt is a focused reading instruction tuned to its category cluster,
including category indicators and asymmetry guidance (e.g.\ liability caps require
recall of monetary thresholds, renewal clauses require recall of conditional triggers).
The specialist returns a verbatim-span answer in the same schema as the baselines. M1
thus combines two architectural commitments: routing and decomposition.

\paragraph{M6 --- Consolidated specialist (ablation).}
M6 is the architectural ablation control. Its prompt is the textual union of the three
specialist prompts of M1, concatenated into a single agent that handles every category
itself, with no router, no orchestrator, and no per-category branching. The category
indicator is selected based on the request, but no decomposition is performed. The point
of M6 is methodological: it lets us decide, in Chapter~5, whether any benefit M1 enjoys
over B1 is attributable to architectural decomposition or merely to the more elaborated
prompt content that the specialists carry. Concretely, M1 vs.\ M6 isolates
``architecture'', and M6 vs.\ B1 isolates ``prompt content''.

\section{Implementation Stack}\label{sec:sys:stack}

\paragraph{Orchestration.}
All configurations are implemented as LangGraph \texttt{StateGraph}s~\cite{langgraph2024},
which let us express the orchestrator/router/specialist topology as typed nodes with
declared transitions. Each node consumes and emits a structured run-state object
containing the contract, the category, intermediate decisions (e.g.\ classifier verdict,
routed specialist), and the in-progress answer. The graph topology is the executable
representation of the architecture diagrams in Section~\ref{sec:sys:patterns}.

\paragraph{Model client.}
A unified multi-provider model client wraps Anthropic, OpenAI, Google Vertex AI, and
Ollama under one synchronous and one streaming interface. The client supports 19 models
in total but only the 11 listed in Section~\ref{sec:meth:models} are used in official
runs. Per call it records prompt tokens, completion tokens, cache-read tokens (where
supported), wall-clock latency, and provider-billed cost. These diagnostics are written
to the trace and aggregated downstream.

\paragraph{Prompt management.}
Prompts are stored in a YAML registry with a stable schema. The registry has two roots,
\texttt{prompts/baselines/} (B1, B4) and \texttt{prompts/specialists/} (the three M1
specialists, the M6 consolidated prompt, the M0 classifier and extractor), and a
separate \texttt{indicators/} directory holding the per-category indicator strings
shared across configurations. The registry is loaded once per run and version-pinned in
the run metadata.

\paragraph{Observability.}
Every model call is traced via Langfuse~\cite{langfuse2024} with a structured trace
format that records the configuration ID, the model identifier, the prompt template
version, the input tokens, the output tokens, the parsed answer, and the routing or
classification decision (where applicable). For M1, the trace records the orchestrator
decision, the routed specialist, and the specialist's answer in a parent-child trace
structure. We measured trace completeness (fraction of model calls for which all
expected trace fields were populated) at $95.2\%$ for M1 across the official runs; the
remaining $4.8\%$ are largely due to provider-side trace truncations on long contracts
and are flagged in the run logs.

\paragraph{Data and persistence.}
The CUAD loader implements the 100{,}000-character filter, the natural-prior sampler,
and the tier partition of Section~\ref{sec:meth:sampling}. It is a deterministic
component keyed by \texttt{seed=42}. Per-run outputs are persisted as JSONL files
containing one record per (sample, configuration, model) prediction with full trace
identifiers, allowing the statistical pipelines (Section~\ref{sec:meth:stats}) to be
re-run from raw logs.

\paragraph{Language and tooling.}
The codebase is Python, managed with \texttt{uv}, and the statistical pipelines are
built on \texttt{pandas}~\cite{pandas2020}, \texttt{numpy}~\cite{harris2020array}, and
\texttt{scipy}~\cite{virtanen2020scipy}. Bootstrap resampling uses \texttt{numpy}
random generators with run-level seeds for reproducibility.

\section{Traceability and Auditability}\label{sec:sys:trace}\label{sec:auditability}

A property of multi-stage agentic systems that is frequently advertised but rarely
characterised is \emph{auditability}: the claim that decomposition produces inspectable
artefacts that enable reviewer trust. We treat auditability as a system property of the
implementation rather than as a hypothesis to be tested, and document here what the
system records and what trade-offs that recording entails. This section addresses
Sub-RQ3 of Chapter~1.

For each clause-level decision, the system captures: (i) the configuration ID and model
identifier, (ii) the prompt template version and the rendered prompt, (iii) the model's
full completion, (iv) the parsed final answer with the parser version, (v) the routing
decision (M1) or classification verdict (M0) where applicable, and (vi) the per-call
token, cost, and latency diagnostics. M1 traces are structured as parent-child
relationships: the orchestrator's routing call is the parent, and the routed
specialist's call is its child. M0 traces analogously expose the classifier as parent
and the extractor (when invoked) as child.

This trace structure supports three concrete auditability properties. First,
\emph{decision attribution}: any final answer can be traced to the specific agent that
produced it and the specific prompt version that elicited it. Second,
\emph{counterfactual replay}: because prompts and model identifiers are pinned per call,
re-running the trace with a different model is mechanical, which lets a reviewer ask
``would this contract have been flagged by a stronger model?'' without re-running the
full pipeline. Third, \emph{rejection forensics}: when M0 silently returns ``no related
clause'' on a category, the classifier's verdict is recorded explicitly, which makes the
failure mode of Section~\ref{sec:res:cascade} (over-rejection) directly visible at
audit time.

The trade-offs are substantive. Verbose traces increase storage cost (the JSONL
persistence layer is on the order of gigabytes per full sweep) and, when
chains-of-thought are recorded, can themselves contain sensitive contract content; we
therefore distinguish between trace fields that are safe to surface to operators
(decisions, token counts) and those that should be access-controlled (full prompts and
completions). Auditability is also a property of the \emph{system}, not of the
underlying model: the trace is faithful to the call graph but does not certify that the
model's stated reasoning is its actual reasoning.

\section{Experimental Infrastructure}\label{sec:sys:infra}

\paragraph{Loader and sampling.}
The CUAD loader is the single source of truth for which contracts and which categories
enter a run. It enforces the 100{,}000-character filter, applies the natural-prior
sampler, and emits the tier label for every sample. For the official runs the loader
produces approximately 2{,}999 samples per run and is identical across configurations,
ensuring that pairing in Chapter~5's statistical pipeline is well-defined.

\paragraph{Statistical aggregation.}
A separate aggregation pipeline reads the per-run JSONL outputs, computes the
per-sample correctness vectors, and emits the McNemar contingency tables, the paired
bootstrap distributions, and the BH-corrected $p$-value tables. Aggregation is
deterministic given the JSONL inputs and a fixed bootstrap seed. The pipeline emits
machine-readable summaries that feed both the thesis tables and the dashboard.

\paragraph{Dashboard.}
A Next.js dashboard is provided for interactive exploration of experiment outputs,
including per-model $F_2$ comparisons, per-tier breakdowns, per-call cost and latency
views, and trace inspection by sample ID. The dashboard is not part of the statistical
pipeline but is the primary interface used during error analysis and was the channel
through which the over-extraction failure mode (Chapter~5) was first surfaced visually.

\paragraph{Test suite.}
An 81-test suite covers the metric implementations ($F_1$/$F_2$, Jaccard, laziness,
over-extraction), the CUAD loader (filter, sampling determinism, tier assignment), the
agent topologies (B1, B4, M0, M1, M6 each have at least one end-to-end test against a
stub model), and the statistical pipelines (McNemar exact, paired bootstrap, BH
correction, intersection logic for partial runs). The full suite passes prior to every
official run. Notebook artefacts
(\texttt{01\_data\_exploration.ipynb},
\texttt{02\_workflow\_test.ipynb},
\texttt{03\_baseline\_calibration.ipynb},
\texttt{04\_multiagent\_experiment.ipynb})
are retained as reproducible exploratory entry points.

\paragraph{Reproducibility.}
Reproducibility is supported by three commitments: deterministic sampling under
\texttt{seed=42}; explicit version pinning of model identifiers in the run metadata;
and persistence of the full trace bundle (prompts, completions, parsed answers,
diagnostics) so that downstream statistics can be regenerated from raw logs without
rerunning models.
